% frontmatter/polish/mapa-del-curso.tex -- "Plan roku": el mapa del curso
% de frontmatter/mapa-del-curso.tex, en polaco.
\section*{Plan roku}

\begin{center}
\renewcommand{\arraystretch}{1.4}
\begin{tabularx}{\linewidth}{@{}c c c c X@{}}
\toprule
\textbf{Część} & \textbf{Dni} & \textbf{Pora roku} & \textbf{Liczby} & \textbf{Co się dzieje} \\
\midrule
1 & 1--65 & jesień & od 0 do 10 & Co tydzień nowa liczba: pisana po
śladzie, liczona, kolorowana i szukana. Lucía ma na torcie siedem
świeczek, a część kończy się Wigilią. \\
2 & 66--130 & zima & do 20 & Dziesięć i liczby od 11 do 20; więcej,
mniej i tyle samo; liczba przed i liczba po; pierwsze dodawanie, przez
łączenie grup na obrazkach. \\
3 & 131--195 & wiosna & do 50 & Dodawanie i odejmowanie do 10; dzielenie
liczby na dwie części; liczenie dwójkami, piątkami i dziesiątkami. \\
4 & 196--260 & lato & do 100 & Dziesiątki do 100; dodawanie i
odejmowanie do 20; pieniądze, zegar i mierzenie, na wsi u babci Rosy. \\
\bottomrule
\end{tabularx}
\end{center}

\vspace{6mm}
Każda część ma 13 tygodni po 5 dni -- także w czasie wakacji: zeszyt to
nadal jedna strona dziennie, od poniedziałku do piątku, w Boże
Narodzenie i latem. Nowa liczba nie przychodzi, dopóki poprzednia nie
była przez cały tydzień pisana po śladzie, liczona, kolorowana i
szukana.
\newpage
